\section{Proposed Approach: Meshtastic Parallel Network}\label{sec:our_approach}
\subsection{Technical Architecture}
We propose \textbf{Meshtastic}, a decentralized mesh network utilizing the \textbf{LoRa physical layer} and Chirp Spread Spectrum (CSS) \cite{meshtastic_overview}. Signals consist of frequency-sweeping ``chirps'' with a 16-up-chirp Preamble for synchronization and a specific Sync Word (0x2B) to distinguish the network \cite{meshtastic_overview}.

\subsection{Decentralized Metadata Protection}
The primary advantage is the complete bypass of Cellular Operator Data Systems \cite{meshtastic_overview}. In a mesh architecture, there is no central database logging ``Top 2 BS'' anchors or hourly traffic windows the features required by SURE and MoCo for re-identification \cite{duan2025sure, duan2025moco, meshtastic_overview}. The SURE framework relies on extracting BS ID sequences, association durations, and traffic levels from centralized operator BSS/OSS databases \cite{duan2025sure}. Similarly, MoCo requires access to cellular signaling traces of the form $\langle$user ID, BS ID, timestamp$\rangle$ stored by operators \cite{duan2025moco}. By removing the centralized collection point entirely, Meshtastic eliminates the cellular operator metadata source that makes these attacks possible.

\subsection{Encryption and Routing}
Meshtastic payloads are encrypted using \textbf{AES-256-CTR encryption} over LoRa, with a different key for each channel \cite{meshtastic_overview}. Only nodes with the correct channel key can interpret data, while hop limits constrain message propagation, reducing the volume of accessible metadata compared to centralized operator logging \cite{meshtastic_overview}.

\subsection{Justification of Selection}
The choice of a parallel radio network over software-side mitigations is motivated by the fundamental limitations of the latter. Approaches such as noise injection and differential privacy attempt to obscure behavioral patterns at the data level; however, they suffer from an inherent trade-off: sufficient noise to defeat re-identification may render the data unusable for its intended purpose. The SURE paper demonstrates that even with k-anonymity (identifier suppression) applied, AUC scores remain above 0.9, indicating that anonymization alone is insufficient \cite{duan2025sure}. Moreover, techniques like ASTRA \cite{chhipa2025astra} and RoCL \cite{kim2020rocl} harden the model against adversarial perturbations but do not address the root cause: \textbf{the centralized logging of spatio-temporal metadata by cellular operators}. Meshtastic addresses this root cause by operating on an entirely separate radio infrastructure, ensuring that no cellular operator ever collects the behavioral traces in the first place.

\section{SWOT Analysis}\label{sec:swot}

To rigorously evaluate the position of the proposed Meshtastic parallel network against existing cellular infrastructure, we conduct a SWOT (Strengths, Weaknesses, Opportunities, Threats) analysis. This comparative framework establishes the credibility of our proposal by systematically examining both internal attributes and external factors.

\begin{table}[htbp]
\caption{SWOT Analysis: Meshtastic Parallel Network vs. Cellular}
\begin{center}
\small
\begin{tabular}{@{}lp{3cm}p{3cm}@{}}
\toprule
 & \textbf{Internal} & \textbf{External} \\
\midrule
\textbf{S/O} & No centralized operator logs; metadata privacy; AES-256-CTR encryption \cite{meshtastic_overview}. & Growing demand for off-grid, high-privacy communication solutions. \\
\midrule
\textbf{W/T} & Lower data rates; limited LoRa range; no PFS \cite{meshtastic_overview}. & Signal triangulation; RF fingerprinting; LoRa band jamming. \\
\bottomrule
\end{tabular}
\end{center}
\end{table}

\textbf{Strengths.} The most significant internal advantage is the complete absence of a centralized operator data system. As demonstrated by SURE \cite{duan2025sure}, re-identification attacks require access to BS association sequences and traffic usage patterns stored in BSS/OSS databases. By removing this centralized collection point, Meshtastic eliminates the ``Top 2 BS'' spatial anchors and hourly traffic windows that enable behavioral fingerprinting. Additionally, AES-256-CTR shared-key encryption ensures payload confidentiality at the application layer \cite{meshtastic_overview}.

\textbf{Weaknesses.} LoRa operates in the sub-GHz ISM bands with data rates limited to approximately 5.47 kbps at the highest spreading factor, making it unsuitable for high-bandwidth applications such as video streaming or large file transfers \cite{meshtastic_overview}. The propagation range, while superior to WiFi in line-of-sight conditions, is constrained in dense urban environments with significant signal attenuation. Furthermore, the mesh topology introduces latency with each additional hop. From a security standpoint, Meshtastic does not provide Perfect Forward Secrecy, message integrity verification, or sender authentication \cite{meshtastic_overview}, meaning that an attacker with access to the channel key can impersonate any node on the network and decrypt passively collected traffic. In the threat model addressed by this report, however, these cryptographic deficiencies are secondary concerns: the primary objective is to eliminate the centralized cellular metadata stream that enables mass behavioral profiling. An attacker who compromises a Meshtastic channel key can surveil at most the traffic of that single mesh group, whereas a compromised cellular operator database exposes the behavioral records of millions of users.

\textbf{Opportunities.} The increasing demand for operational security in privacy-sensitive domains creates a growing market for off-grid communication solutions. Recent data breaches at major operators (AT\&T, T-Mobile, US Cellular) \cite{duan2025sure} have highlighted the vulnerability of centralized data repositories, strengthening the case for decentralized alternatives. Application domains such as journalism, activism, and disaster relief stand to benefit from infrastructure-independent communication.

\textbf{Threats.} While Meshtastic avoids cellular metadata collection, it remains potentially vulnerable to RF-level attacks. Signal triangulation using directional antennas could determine the physical location of transmitting nodes. Future advances in RF fingerprinting which identify unique hardware characteristics of individual radio transmitters could theoretically enable device-level identification without requiring centralized metadata logs. Additionally, intentional signal jamming in the LoRa frequency bands could disrupt mesh network operations.
